\documentclass[UTF8,AutoFakeBold]{szu-master}
\renewcommand{\titleCN}{基于XXX研究}
\renewcommand{\titleEN}{The Research Base on XXXX}

% 源代码地址：https://github.com/Cao-Wuhui/SZU_Latex_Master_Template

% 文本
\usepackage{ctex}
\usepackage{fancyhdr}
%可跨页表格引入
\usepackage{longtable}
\usepackage{lastpage}
% dgq package
\usepackage{makecell}

\usepackage{fontspec}
%圆
%\usepackage{xeCJK}

%%adding package 边距设置
\usepackage[left=31.7mm,right=31.7mm,top=25.4mm,bottom=25.4mm]{geometry}

\usepackage{titlesec}
\usepackage{setspace} %使用调整行距的包
\usepackage{subfigure,floatrow}
\usepackage{url}
\usepackage{caption}
\usepackage{subcaption}
\usepackage{caption,cite,amsbsy,amsmath,amsfonts,multirow,color,array,booktabs,colortbl}
\usepackage{algorithm}
\usepackage{algorithmic}
\usepackage{pdfpages}
\usepackage{changepage}

\usepackage{titlesec}

% \usepackage{algorithmicx,algorithm}
\usepackage{epsfig,bm}
\usepackage{slashbox}
\usepackage{enumitem}
\usepackage{amssymb,amsthm,amssymb}
\usepackage{graphicx,subfigure}
\usepackage{graphicx}
\usepackage{ulem}
\usepackage{floatrow}
\usepackage{arydshln} % 加载arydshln宏包
\usepackage{tabularx}
\usepackage{chngcntr}
\floatsetup[table]{capposition=top}
\floatsetup[figure]{capposition=bottom}
\newfloatcommand{capbtabboxTable}{table}[][\FBwidth]
\newfloatcommand{capbtabboxFigure}{figure}[][\FBwidth]
\newcommand{\norm}[1]{\left\lVert #1 \right\rVert}

% 脚注样式 ①②③
% 自动检测并选择圈号覆盖更全的符号字体（按优先级回退）
\IfFontExistsTF{Noto Sans Symbols 2}
  {\newfontfamily\circnumfont{Noto Sans Symbols 2}}
  {\IfFontExistsTF{Segoe UI Symbol}
    {\newfontfamily\circnumfont{Segoe UI Symbol}}
    {\IfFontExistsTF{Symbola}
      {\newfontfamily\circnumfont{Symbola}}
      {\IfFontExistsTF{Arial Unicode MS}
        {\newfontfamily\circnumfont{Arial Unicode MS}}
        {\newfontfamily\circnumfont{DejaVu Sans}}
      }%
    }%
  }
  
\newcommand{\circledfootnotenum}[1]{%
  {\circnumfont
    \ifcase#1
      \or ①\or ②\or ③\or ④\or ⑤\or ⑥\or ⑦\or ⑧\or ⑨\or ⑩%
      \or ⑪\or ⑫\or ⑬\or ⑭\or ⑮\or ⑯\or ⑰\or ⑱\or ⑲\or ⑳%
      \or ㉑\or ㉒\or ㉓\or ㉔\or ㉕\or ㉖\or ㉗\or ㉘\or ㉙\or ㉚%
      \or ㉛\or ㉜\or ㉝\or ㉞\or ㉟\or ㊱\or ㊲\or ㊳\or ㊴\or ㊵%
      \or ㊶\or ㊷\or ㊸\or ㊹\or ㊺\or ㊻\or ㊼\or ㊽\or ㊾\or ㊿%
    \else
      #1%
    \fi
  }%
}
\renewcommand{\thefootnote}{\circledfootnotenum{\value{footnote}}}

\newtheorem{definition}{定义}[] % 跟随三级标题排序
\newtheorem{theorem}{定理}[]  %跟随二级标题排序
\newtheorem{lemma}[theorem]{引理}
\newtheorem{corollary}[theorem]{推论}
\newtheorem{proposition}{命题}[]


% 重新调整参考文献项目之间的行距
\usepackage{bibspacing}
% \setlength{\bibspacing}{1\baselineskip}

%给目录和参考文献加上超链接
%\usepackage[hidelinks]{hyperref}

%引用
\newcommand{\upcite}[1]{{\textsuperscript{\cite{#1}}}}

%defines the height of the header
% \setlength{\headheight}{15pt}
\setlength{\headheight}{23pt}
\setlength{\headsep}{7.25mm}


\hypersetup{
	pdfborder = {0 0 0} 
}

%\titleformat{\section}{\center\zihao{3}\heiti\bfseries\sectionef}{第\,\zhnumber{\thesection}\,章}{1em}{}

% 在导言区定义新字体族
% \newfontfamily\sectionef{Arial}

\ctexset{
		section={
				format+ = \heiti\bfseries\zihao{3}\sectionef,
				name = {第,章},
				number = \chinese{section}
	}
}	


\titleformat{\subsection}{\zihao{-3}\heiti\bfseries\sectionef}{\thesubsection}{1em}{}

\titleformat{\subsubsection}{\zihao{4}\bfseries\fontspec{songti-Bold.otf}}{\thesubsubsection}{1em}{}


% 定义目录标题样式
\titlecontents{section}[0em]{\bfseries\fontspec{songti-Bold.otf}\zihao{-4}}{\contentspush{\thecontentslabel\hspace{1em}}}{}{\titlerule*[0.4pc]{.}\contentspage\vspace{1.6pt}}%

%\titlecontents{section}[0em]{\bfseries\zihao{-4}}{\contentspush{\thecontentslabel\hspace{1em}}}{}{\titlerule*[0.4pc]{.}\contentspage\vspace{1.6pt}}%

\titlecontents{subsection}[1em]{\fontspec{simsun.ttc}\zihao{-4}}{{\thecontentslabel}\hspace{1em}}{}{\titlerule*[0.4pc]{.}\contentspage\vspace{1.6pt}}

\titlecontents{subsubsection}[1.8em]{\fontspec{simsun.ttc}\zihao{-4}}{\thecontentslabel\hspace{1em}}{}{\titlerule*[0.4pc]{.}\contentspage\vspace{1.6pt}}




%document begin
\begin{document}
	
%封面 自行配置生成pdf文件导入即可
\includepdf[pages={1,2}]{chapter/masterPage.pdf}
	
% 全局页眉页脚：确保摘要等多页内容的后续页保持 fancy 样式
% 解决 issue: https://github.com/Ftine/SZU_Latex_Master_Template/issues/19
\pagestyle{fancy}
\fancyhead[C]{\songti\zihao{5}{\titleCN}}
\fancyhead[C]{\songti\zihao{5}{\titleEN}}
	
%中英文摘要
\input{chapter/abstract.tex} 
	
\setlength{\baselineskip}{23pt} %设置行距23磅
	
\vspace*{-1\baselineskip}
\tableofcontents	%目录

%缩略语
\input{chapter/sign_and_ABC.tex}

%此处为原cls文件中有关后续正文页面的页眉页脚设置
\newpage
\pagenumbering{arabic}
\pagestyle{fancy} % 页眉横线样式
\fancyhead{}

\fancyhead[C]{\songti\zihao{5}{\titleCN}}
% 全局页眉 设置为章节标题
\fancyhead[C]{\songti\zihao{5}\nouppercase\leftmark}
%全局页号设置
\fancyfoot[C]{\songti\zihao{5}{\thepage}}
\setcounter{page}{1}

%%%%%%%%% 图索引和表索引
% \renewcommand\listfigurename{\songti\zihao{2}{\textbf{图~~目~~录}}}
% \renewcommand\listtablename{\songti\zihao{2}{\textbf{表~~目~~录}}}
	
% \renewcommand{\figurename}{图}
% \renewcommand{\tablename}{表}
	
\makeatletter %
\@addtoreset{equation}{section}
\makeatother  %
% 公式更改为3-2样式
\renewcommand\theequation{{\thesection}-{\arabic{equation}}}
	
%%%%%%%%%% 表每节重新标号
\makeatletter %
\@addtoreset{table}{section}
\makeatother  %
\renewcommand{\thetable}{\thesection-\arabic{table}}
\captionsetup[table]{labelsep=space, labelfont=bf}
%%%%%%%% 图每节重新标号
\makeatletter %
\@addtoreset{figure}{section}
\makeatother  %
\renewcommand{\thefigure}{\thesection-\arabic{figure}}
\captionsetup[figure]{labelsep=space, labelfont=bf}
%labelfont=normalfont}
% subfigure重新编号\ref
\makeatletter
\renewcommand{\p@subfigure}{\thefigure}
\makeatother


%adding chapters 1-5 从1开始计数
\setcounter{page}{1}

\input{chapter/cp1.tex}
\input{chapter/cp2.tex}
\input{chapter/cp3.tex}
\input{chapter/cp4.tex}
\input{chapter/cp5.tex}


\clearpage
%参考文献
\fancyhead{}
\fancyhead[C]{\songti\zihao{5}{参考文献}}
\fancyfoot[C]{\songti\zihao{5}{\thepage}}

\begin{spacing}{1.8}
	\bibliographystyle{gbt7714-2005} 
	\vspace*{-1\baselineskip}
	\bibliography{ref_paper}
\end{spacing}



\clearpage
\fancyhead{}
\fancyhead[C]{\songti\zihao{5}{致谢}}
\fancyfoot[C]{\songti\zihao{5}{\thepage}}
\input{chapter/thanks.tex}

\clearpage
\fancyhead{}
\fancyhead[C]{\songti\zihao{5}{攻读硕士学位期间的学术成果}}
\fancyfoot[C]{\songti\zihao{5}{\thepage}}
\input{chapter/myPapers.tex}

% 添加学术评语目录条目
%\addtocontents{toc}{\protect\vspace{1.6pt}}
%\addtocontents{toc}{\noindent\protect\fontspec{songti-Bold.otf}\protect\zihao{-4}\protect\normalfont 附：指导教师对研究生学位论文的学术评语\par}
% 添加答辩决议书目录条目
%\addtocontents{toc}{\protect\vspace{1.6pt}}
%\addtocontents{toc}{\noindent\protect\fontspec{songti-Bold.otf}\protect\zihao{-4}\protect\normalfont ~~~~~~~~答辩委员会决议书\par}
% 添加学术评语页面
%\clearpage
%\includepdf[pages={1}]{chapter/comment.pdf}
% 添加答辩决议书页面
%\clearpage
%\includepdf[pages={1,2}]{chapter/decision.pdf}

%当需要在论文末尾插入相关的文件，可以生成pdf文件格式进行附加，附加一页就是1，多页就是1，2，3
%\includepdf[pages={1}]{chapter/comment.pdf}
%\includepdf[pages={1}]{chapter/decision.pdf}

\end{document}

% 源代码地址：https://github.com/Ftine/SZU_Latex_Master_Template
